\section{Discussion}\label{sec:disc}

\subsection{Observations}

The surprise is not that abstracts compress; it is the asymmetry.
Abstracts under-report methodology most (benchmarking flags: 5 in
abstracts, 20 in bodies), and full-text binary coding over-reports in
exactly the mirror image.
Cheap reads and naive deep reads are both wrong, in opposite
directions; only the thresholded middle agrees with a human coder.
This is why rite's gate is stated on the thresholded coding, not on
either extreme.

Second, the pipeline's byproducts kept becoming the argument.
The download rate was meant as logging and became Finding 1; the
saturation artifact was a bug and became Finding 3.
A reproducible search leaves evidence lying around, and evidence is
what papers are made of.

Third, the pipeline's most consequential byproduct may be the problem
named in the FAQ (Section~\ref{sec:faq}): title-plus-abstract
miscoding touches indexing, reading decisions, and reviewer bidding,
and SE has barely begun to discuss it~\cite{brereton07}.

\subsection{Threats to validity}

\textbf{Instrument.}
The 0.5-per-1000-words cutoff is ours, unvalidated beyond this corpus;
the promised sensitivity analysis has not yet run.
Keyword coding is a draft that the method itself says to hand-audit.

\textbf{Selection.}
One goal, one field, 22 full texts; the paywall excluded most
classics, so Finding 2 may differ on an older, closed corpus.

\textbf{Author as instrument.}
The loop's expensive steps (10, 12) were run by the person who built
the loop.
This paper is deliberately pre-critics: it is paper[1] of its own
Table~\ref{tab:loop}, awaiting steps 10--13.

\subsection{Future work}

Run this paper through its own critics (steps 10--12), publish the
critic reports and the auto-corrected paper[2] beside this draft;
snowball forward and backward from the anchor review to complete the
related work; run the threshold sensitivity analysis; add the
five-year self-coding exhibit.
